\section{Fundamentos e trabalhos relacionados}

\pendente{Escopo desta seção fixado pela orientação em 19/08/2026:
administração pública, prefeituras e licitação. O critério dado foi escrever o
necessário para o leitor entender o que está sendo entregue, e nada além. O
conteúdo de computação --- recuperação de informação, representação vetorial e
trilha de auditoria --- não entra aqui: é contribuição, e vive na
\hyperref[sec:metodologia]{metodologia}.}

\pendente{Cada subseção nasce de um fichamento em \code{TCC/papers/} e de uma
nota conceitual em \code{TCC/referência/}. Não escrever aqui antes de ler a
fonte --- a lição registrada na IC é que reconstruir metadados de memória gera
erro de citação.}

\subsection{Contratações públicas e a fase preparatória}

\pendente{Lei nº 14.133/2021: fase preparatória, ETP, TR e valor estimado.
Fonte primária: o texto da lei, já conferido no portal oficial. Fontes
secundárias: manuais e orientações de órgãos de controle.}

\subsection{Pesquisa de preços em contratações públicas}

\pendente{Parâmetros admitidos para a pesquisa de preços e tratamento
estatístico dos valores obtidos. Fonte primária: a Instrução Normativa
SEGES/ME nº 65/2021, já conferida no portal oficial.}

\subsection{Trabalhos relacionados}
\label{subsec:trabalhos-relacionados}

O quadro a seguir é comparativo e não exaustivo. A restrição de extensão não
acomoda protocolo de revisão sistemática com bases, expressão de busca,
critérios de inclusão e janela temporal, e o produto de um protocolo desses não
caberia em duas páginas. O protocolo e o resultado bruto do levantamento ficam
registrados fora do texto principal, e esta subseção apresenta os trabalhos que
delimitam o espaço ocupado aqui.

\pendente{O quadro abaixo traz quatro trabalhos, e não os cinco a oito
pretendidos. Todos foram lidos integralmente e tiveram os metadados conferidos
na fonte primária. Outros sete estão conferidos e fichados, e entram nas
próximas semanas --- a limitação é deliberada, para que cada trabalho citado
possa ser discutido em profundidade.}

\begin{table}[H]
  \centering
  \caption{Trabalhos relacionados e posição deste trabalho}
  \label{tab:trabalhos-relacionados}
  \begin{ufopaespacosimples}
  \footnotesize
  \begin{tabularx}{\textwidth}{@{}lXXl@{}}
    \toprule
    \textbf{Trabalho} & \textbf{Etapa atendida} & \textbf{Tratamento de preço} &
      \textbf{País} \\
    \midrule
    \cite{siciliani2023apoio} & Ciclo inteiro; consulta a editais e contratos &
      Não trata & Itália \\
    \cite{brum2024agrupamento} & Agrupamento de itens de compra pública &
      Não trata preço; agrupa descrições & Brasil \\
    \cite{faria2025precos} & Fase preparatória; cálculo do valor estimado &
      Tratamento estatístico sobre cotações já reunidas & Brasil \\
    \cite{vandermeer2026legalcheck} & Fora de compras; parecer jurídico
      municipal & Não trata & Países Baixos \\
    \midrule
    \textbf{Este trabalho} & \textbf{Fase preparatória; peças e triagem de
      preço} & \textbf{Recuperação por significado sob filtro estruturado} &
      \textbf{Brasil} \\
    \bottomrule
  \end{tabularx}
  \end{ufopaespacosimples}
  \fonte{Elaborado pelo autor.}
\end{table}

O sistema de apoio à decisão descrito por \citaautor{siciliani2023apoio},
desenvolvido com uma região italiana, é o precedente mais direto do emprego de
busca por significado em compras públicas: reúne dados estruturados da base
nacional de contratos e documentos de editais, e oferece busca que combina
palavra e proximidade semântica, com sugestão de conceitos relacionados para
cobrir o vocabulário que o usuário não domina. O alvo, porém, é a consulta a
processos e empresas, não a composição do valor estimado, e o preço não é
tratado.

\citaautor{brum2024agrupamento} atacam o problema de representação que está no
centro deste trabalho, ainda que com outro objetivo. Para agrupar itens de
compra pública que são o mesmo produto escrito de formas diferentes, comparam
oito representações de texto sobre mais de dois milhões de descrições em
português, e propõem uma variante que pondera os componentes que definem um
item: substantivo, unidade de medida e número. Dois achados deles orientam o
desenho adotado aqui. O primeiro é que o sinal lexical é forte nesse tipo de
descrição, a ponto de uma heurística simples organizar a maior parte do conjunto:
sem ela, de 43\% a 45\% dos itens ficam sem grupo. O segundo é que método simples
fica próximo de modelo sofisticado, o que desaconselha fixar a representação
antes de medir. A diferença de objetivo importa: agrupar é
operação sobre o acervo inteiro, enquanto este trabalho precisa recuperar
candidatos para um item específico, sob filtro, e com decisão humana ao fim.

\citaautor{faria2025precos} são o correlato brasileiro mais próximo quanto ao
preço. Desenvolvem software que calcula o preço de referência a partir das
cotações já coletadas, criticando os três critérios usuais e selecionando entre
tratamentos estatísticos alternativos. A diferença de escopo é o ponto: o
trabalho deles começa onde as cotações já estão reunidas, e é exatamente a etapa
anterior --- decidir, item por item, quais preços são comparáveis ao item
demandado --- que concentra o esforço medido aqui.

Há ainda uma divergência de critério que precisa ser registrada, porque afeta a
avaliação. Aquele trabalho mede sucesso pela proximidade entre o preço de
referência calculado e o preço efetivamente homologado. No setor estudado, o
valor estimado funciona como teto, e a servidora preserva margem acima do preço
praticado de propósito. Adotar aquela métrica mediria este sistema contra um
alvo que a prática do domínio considera indesejável, e por isso ela não é
adotada.

\citaautor{vandermeer2026legalcheck} descrevem sistema em uso no departamento
jurídico de uma prefeitura europeia, que recupera pareceres anteriores e
legislação, gera minuta e a submete a revisão obrigatória de um jurista, que
permanece autor do documento. É o precedente mais próximo quanto à geração
assistida de peça administrativa em prefeitura, inclusive por trazer avaliação em
uso real por seis meses, com medida de tempo e de aceitação. O domínio é outro
--- resposta a objeção de cidadão, não contratação ---, e não há tratamento de
preço, mas a arquitetura de revisão obrigatória com autoria humana preservada é
a mesma adotada aqui.

\subsection{Sistemas em operação}
\label{subsec:sistemas-em-operacao}

Trabalho acadêmico não é o único ponto de comparação, e para um sistema que se
propõe a apoiar uma rotina administrativa talvez nem seja o principal. Esta
subseção examina ferramentas em uso, de governo e comerciais, cuja fonte é
documentação do mantenedor, norma, tutorial oficial ou observação direta da
interface pública --- e não literatura.

\begin{table}[H]
  \centering
  \caption{Ferramentas em operação e o mecanismo de busca de cada uma}
  \label{tab:sistemas-em-operacao}
  \begin{ufopaespacosimples}
  \footnotesize
  \begin{tabularx}{\textwidth}{@{}XXXl@{}}
    \toprule
    \textbf{Ferramenta} & \textbf{Etapa atendida} & \textbf{Como encontra o
      item} & \textbf{Acesso} \\
    \midrule
    Módulo federal de estudo técnico preliminar & Redação da peça, por
      formulário & Código de catálogo & Órgãos federais \\
    Módulo federal de termo de referência & Redação da peça, por formulário &
      Código de catálogo & Órgãos federais \\
    Assistente federal de normas & Dúvida sobre legislação de contratações &
      Recuperação sobre normas, não sobre item & Órgãos federais \\
    Pesquisa de preços federal & Composição do valor estimado & Texto casado por
      caracteres sobre o catálogo & Público e com login \\
    Portal nacional de contratações & Consulta a contratações publicadas & Texto
      por palavra; devolve a contratação, não o item & Público \\
    Bancos de preços comerciais & Coleta, cálculo e relatório & Texto por
      palavra & Assinatura \\
    Geradores comerciais de peça & Redação com modelo de linguagem & Não
      declarado & Assinatura \\
    \midrule
    \textbf{Este trabalho} & \textbf{Triagem de preço e redação das peças} &
      \textbf{Significado, sob filtro estruturado} & --- \\
    \bottomrule
  \end{tabularx}
  \end{ufopaespacosimples}
  \fonte{Elaborado pelo autor a partir de documentação oficial, tutoriais dos
    mantenedores e observação das interfaces públicas em outubro de 2026.}
\end{table}

Três constatações decorrem desse levantamento, e duas delas limitam o que este
trabalho pode reivindicar.

Gerar as peças da fase preparatória com modelo de linguagem já é produto de
prateleira. Há ao menos oito produtos comerciais que o anunciam. Nenhum deles
publica método, modelo, base de recuperação ou avaliação, mas a ausência de
publicação não torna a função inédita. A contribuição deste trabalho, portanto,
não está em gerar o documento, e sim em que cada afirmação do documento gerado
tenha origem declarada e conferível.

Filtro estruturado sobre item também já existe nas ferramentas federais, que
filtram por unidade de medida, quantidade, localização e data. O diferencial não
é o filtro: é combiná-lo com recuperação por significado e registrar a decisão de
quem escolheu.

A terceira constatação abre espaço em vez de fechá-lo. Nenhuma das ferramentas
consultadas documenta guardar o candidato descartado com o motivo do descarte,
nem reconstituir uma pesquisa tal como ela era em uma data. A norma federal de
pesquisa de preços exige justificar a desconsideração de valores
\cite{brasil2021in65}, e nenhuma ferramenta examinada descreve onde essa
justificativa fica registrada.

\paragraph{A lacuna ocupada.} Dos trabalhos e sistemas acima, nenhum apoia a
decisão, item por item, de quem monta a pesquisa de preços. Os correlatos
brasileiros de preço ou partem das cotações já reunidas, ou analisam sobrepreço
depois da contratação; o precedente de busca semântica em compras públicas não
trata preço; o de geração de peça em prefeitura atua em outro domínio
documental; e nenhuma ferramenta em operação busca o item por significado, pois
todas casam caracteres, código de catálogo ou código de barras. Soma-se a isso
que nenhum deles dispõe de gabarito das escolhas reais de quem cotou, que este
trabalho possui.

\pendente{Esta afirmação de lacuna vale para os quatro trabalhos citados, os sete
já fichados e as ferramentas examinadas. Ela não é afirmação de inexistência:
dez páginas do portal federal de compras e os dois bancos comerciais exigem
acesso autenticado e não puderam ser verificados, e um tutorial oficial declara
conteúdo retirado do ar por período eleitoral até 25/10/2026. Repetir a consulta
depois dessa data.}
